\documentclass[11pt]{article}

\usepackage[margin=1in]{geometry}
\usepackage{amsmath,amssymb,bm}
\usepackage{booktabs}
\usepackage[hidelinks]{hyperref}

\newcommand{\dd}{\mathrm{d}}
\newcommand{\loglike}{\mathcal{L}}
\newcommand{\data}{\mathcal{D}}
\newcommand{\sigmoid}{\operatorname{sigmoid}}
\newcommand{\Normal}{\mathcal{N}}
\newcommand{\Student}{t}

\title{Stage-2 and Stage-3 Independent-Scaling Likelihoods: Blend and Population Modes}
\author{lenscluster implementation note}
\date{\today}

\begin{document}
\maketitle

\section{Purpose}

This note writes the stage-2 and stage-3 likelihoods used for independent/free
scaling galaxies in the two implemented modes:
\renewcommand{\labelenumi}{(\roman{enumi})}
\begin{enumerate}
    \item \textbf{blend}: the existing continuous branch-weight model;
    \item \textbf{population}: the explicit approximate scaling/free
    population-mixture model.
\end{enumerate}
Both modes use the same underlying source-position likelihood.
They differ only in how the independent/free candidate galaxies enter the
lens model and in whether the free probability appears explicitly as a
mixture weight in the likelihood.

Stage 2 is a joint source-plane stage. It always uses the source-plane core
likelihood. Stage 3 is the optional image-plane-preconditioning stage and can
use either source-plane or local-Jacobian mode. In the common sequential setup,
stage 3 uses local-Jacobian mode.

\section{Where stage 2 and stage 3 enter the workflow}

In sequential fitting, the relevant stages are:
\begin{center}
\begin{tabular}{llll}
\toprule
Stage & Run directory & Lens parameters & Core likelihood\\
\midrule
1 & \texttt{stage1\_large\_only} & large-scale halos only & source-plane\\
2 & \texttt{stage2\_joint} & full joint model & source-plane\\
3 & \texttt{stage3\_image\_plane} & full joint model & source-plane or local-Jacobian\\
4 & image-plane stage & full joint model & image-plane\\
\bottomrule
\end{tabular}
\end{center}

Stage 2 is the first full joint stage. It includes the large-scale components,
the scaling-relation member galaxies, and any selected independent/free
candidate branches. Since stage 2 is source-plane only, its independent-scaling
likelihood is
\[
    \ell_{\rm stage2}(\Theta)
    =
    \ell_{\rm src}(\Theta).
\]
The command-line option
\[
    \texttt{--stage2-independent-scaling-likelihood}
    \in \{\texttt{blend},\texttt{population}\}
\]
chooses how the p\_free branches enter this source-plane likelihood.

Stage 3 uses the same full joint parameterization and the same p\_free gate.
The command-line option
\[
    \texttt{--stage3-independent-scaling-likelihood}
    \in \{\texttt{blend},\texttt{population}\}
\]
chooses the p\_free mode for stage 3. If stage 3 is local-Jacobian, the
population mode has to recompute the one-candidate replacement likelihoods
with both source coordinates and local Jacobian covariance terms. If stage 3
is source-plane, it has the same core likelihood as stage 2.

Thus the stage-2 source-plane population mode is mathematically the same
population mixture described below, with
\[
    \ell_{\rm core}=\ell_{\rm src}.
\]
It is usually much cheaper than stage-3 local-Jacobian population mode because
candidate replacements only need source-coordinate likelihoods; they do not
need per-candidate Jacobian replacements and full local covariance
recalculations.

\section{Data and lens mapping}

Let observed image positions be
\[
    \bm{x}_j = (x_j, y_j), \qquad j = 1,\ldots,N_{\rm img}.
\]
Each image belongs to a multiply-imaged family
\[
    f_j \in \{1,\ldots,N_{\rm fam}\}.
\]
The model parameters are denoted collectively by $\Theta$. For a given lens
state, the source-plane position of image $j$ is
\[
    \bm{\beta}_j(\Theta) =
    \bm{x}_j - \bm{\alpha}(\bm{x}_j;\Theta),
\]
where $\bm{\alpha}$ is the total lens deflection field.

For local-Jacobian mode we also use the image-to-source Jacobian
\[
    A_j(\Theta) =
    \frac{\partial \bm{\beta}}{\partial \bm{x}}(\bm{x}_j;\Theta)
    =
    \begin{pmatrix}
    a_{00,j} & a_{01,j}\\
    a_{10,j} & a_{11,j}
    \end{pmatrix}.
\]

Each image has an observed positional scale $\sigma_j$, a reliability
$\rho_j \in (0,1)$, and a Boolean mask $m_j$ indicating whether that image is
used as a constraint. The code clips reliabilities away from exactly zero and
one. A broad outlier component is also included at the image-residual level.

\section{The independent/free gate}

For each independent/free candidate galaxy $i$, the gate probability is
\[
    q_i(\Theta)
    =
    p_{\rm free,i}
    =
    \sigmoid\!\left(
        b_{p(i)}
        +
        a\,\eta_i
        +
        u_i
    \right).
\]
Here:
\begin{itemize}
    \item $p(i)$ is the potfile containing galaxy $i$;
    \item $b_{p(i)}$ is the potfile-level gate intercept;
    \item $a$ is the global magnitude slope;
    \item $\eta_i$ is the standardized magnitude feature for galaxy $i$;
    \item $u_i$ is the per-galaxy local logit offset.
\end{itemize}
The prior on these gate parameters controls the prior trend of $q_i$, but the
two likelihood modes use $q_i$ differently.

\section{Core source-plane likelihood}

Both blend and population modes call a common source-position likelihood after
constructing a set of source coordinates $\{\bm{\beta}_j\}$, and, for
local-Jacobian mode, Jacobians $\{A_j\}$.

\subsection{Source-plane covariance and centroid}

In pure source-plane mode, define
\[
    s_{x,j}^2
    =
    \sigma_j^2\,g_j
    +
    \sigma_{\rm src}^2
    +
    v_{x,j}^{\rm sc}
    +
    c_{\rm floor},
    \qquad
    s_{y,j}^2
    =
    \sigma_j^2\,g_j
    +
    \sigma_{\rm src}^2
    +
    v_{y,j}^{\rm sc}
    +
    c_{\rm floor}.
\]
Here $g_j$ is the cached inverse absolute magnification factor used by the
source-plane approximation, $\sigma_{\rm src}$ is the sampled or fixed source
intrinsic scatter, $v_{x,j}^{\rm sc},v_{y,j}^{\rm sc}$ are scaling-scatter
extra variances, and $c_{\rm floor}$ is a covariance floor. The code optionally
floors $g_j$ using the likelihood stabilizer maximum gain.

The source centroid for family $f$ is the reliability-weighted mean
\[
    \widehat{\bm{\beta}}_f
    =
    \frac{
        \sum_{j:f_j=f} w_j \bm{\beta}_j
    }{
        \sum_{j:f_j=f} w_j
    },
    \qquad
    w_j =
    \frac{\rho_j}{\frac{1}{2}(s_{x,j}^2+s_{y,j}^2)}.
\]
The residual is
\[
    \bm{r}_j =
    \bm{\beta}_j - \widehat{\bm{\beta}}_{f_j}.
\]
The implementation may smoothly cap very large residuals for numerical
stability.

\subsection{Source-plane image contribution}

For the Gaussian residual loss, the in-family source-plane log density is
\[
    \ell_{{\rm fam},j}^{\rm src}
    =
    -\frac{1}{2}
    \left[
        \frac{r_{x,j}^2}{s_{x,j}^2}
        +
        \frac{r_{y,j}^2}{s_{y,j}^2}
        +
        \log(2\pi s_{x,j}^2)
        +
        \log(2\pi s_{y,j}^2)
    \right].
\]
When the Student-t stabilizer is selected, the code replaces this Gaussian
term by the corresponding two-dimensional Student-t log density with the same
quadratic form and log determinant.

The broad outlier log density is
\[
    \ell_{{\rm out},j}
    =
    -\frac{1}{2}
    \left[
        \frac{r_{x,j}^2+r_{y,j}^2}{\sigma_{\rm out}^2}
        +
        2\log(2\pi\sigma_{\rm out}^2)
    \right].
\]
The image-level reliability mixture is
\[
    \ell_j^{\rm src}
    =
    \log\!\left[
        \rho_j \exp(\ell_{{\rm fam},j}^{\rm src})
        +
        (1-\rho_j)\exp(\ell_{{\rm out},j})
    \right].
\]
Thus the source-plane log likelihood is
\[
    \ell_{\rm src}(\Theta)
    =
    \sum_{j=1}^{N_{\rm img}}
    m_j\,\ell_j^{\rm src}.
\]

\section{Core local-Jacobian likelihood}

Local-Jacobian mode uses the same source coordinates, but constructs a full
two-dimensional covariance for each image using the local Jacobian. Define
\[
    C_j =
    \sigma_j^2
    \begin{pmatrix}
    a_{00,j}^2+a_{01,j}^2 &
    a_{00,j}a_{10,j}+a_{01,j}a_{11,j}\\
    a_{00,j}a_{10,j}+a_{01,j}a_{11,j} &
    a_{10,j}^2+a_{11,j}^2
    \end{pmatrix}
    +
    \begin{pmatrix}
    \sigma_{\rm src}^2+v_{x,j}^{\rm sc}+c_{\rm floor} & 0\\
    0 & \sigma_{\rm src}^2+v_{y,j}^{\rm sc}+c_{\rm floor}
    \end{pmatrix}.
\]
The implementation also adds a small gain-floor term when requested, jitters
the determinant to remain positive, and uses $C_j^{-1}$ to estimate the
family centroid. The local-Jacobian centroid solves
\[
    \widehat{\bm{\beta}}_f
    =
    \left[
        \sum_{j:f_j=f}\rho_j C_j^{-1}
    \right]^{-1}
    \left[
        \sum_{j:f_j=f}\rho_j C_j^{-1}\bm{\beta}_j
    \right].
\]
The residual is again
\[
    \bm{r}_j =
    \bm{\beta}_j - \widehat{\bm{\beta}}_{f_j}.
\]
For Gaussian residuals,
\[
    \ell_{{\rm fam},j}^{\rm LJ}
    =
    -\frac{1}{2}
    \left[
        \bm{r}_j^{\mathsf T}C_j^{-1}\bm{r}_j
        +
        \log\!\left((2\pi)^2\det C_j\right)
    \right].
\]
The same broad outlier term is mixed in:
\[
    \ell_j^{\rm LJ}
    =
    \log\!\left[
        \rho_j \exp(\ell_{{\rm fam},j}^{\rm LJ})
        +
        (1-\rho_j)\exp(\ell_{{\rm out},j})
    \right],
\]
and
\[
    \ell_{\rm LJ}(\Theta)
    =
    \sum_{j=1}^{N_{\rm img}} m_j\,\ell_j^{\rm LJ}.
\]

For compactness, let
\[
    \ell_{\rm core}(\Theta)
    =
    \begin{cases}
    \ell_{\rm src}(\Theta), & \text{stage 2, or stage-3 source-plane mode},\\
    \ell_{\rm LJ}(\Theta), & \text{stage-3 local-Jacobian mode}.
    \end{cases}
\]

\section{Blend mode}

In blend mode, $q_i=p_{\rm free,i}$ is a continuous lens-field branch weight.
For a candidate galaxy $i$, let $S_i$ denote its scaling-relation dPIE branch
and $F_i$ denote its independent/free dPIE branch. In the packed lens state the
branch weight is
\[
    w(S_i) = 1-q_i,
    \qquad
    w(F_i) = q_i.
\]
Operationally, this weight multiplies the dPIE mass normalization used by the
branch. Because the deflection and lensing Jacobian are linear in that
normalization, this is equivalent to using the blended candidate contribution
\[
    \bm{\alpha}_i^{\rm blend}
    =
    (1-q_i)\bm{\alpha}_{S_i}
    +
    q_i\bm{\alpha}_{F_i},
\]
and similarly
\[
    J_i^{\rm blend}
    =
    (1-q_i)J_{S_i}
    +
    q_iJ_{F_i}.
\]
The total deflection field is therefore
\[
    \bm{\alpha}^{\rm blend}
    =
    \bm{\alpha}_{\rm other}
    +
    \sum_{i=1}^{N_{\rm cand}}
    \left[
        (1-q_i)\bm{\alpha}_{S_i}
        +
        q_i\bm{\alpha}_{F_i}
    \right],
\]
with the analogous expression for $A_j^{\rm blend}$ in local-Jacobian mode.
In stage 2 source-plane mode, only the blended deflection field and resulting
source coordinates are used; the Jacobian expression is not evaluated by the
core likelihood.

The blend-mode likelihood is a single core likelihood evaluation:
\[
    \boxed{
    \log p(\data\mid \Theta,{\rm blend})
    =
    \ell_{\rm core}\!\left(\Theta;\bm{\alpha}^{\rm blend},A^{\rm blend}\right).
    }
\]

The important consequence is that $q_i$ does not appear as
\[
    \log\left[(1-q_i)L_{\rm scaling}+q_iL_{\rm free}\right].
\]
Instead, it changes the lens field continuously before the usual likelihood is
computed. Thus the posterior samples of $q_i$ are constrained by how the
continuous blended lens field improves the fit; they are not discrete
membership probabilities.

\section{Population mode}

\subsection{Exact latent-branch target}

The exact discrete population model would introduce latent branch assignments
\[
    z_i \in \{0,1\},
    \qquad
    z_i=0 \text{ for scaling},\quad z_i=1 \text{ for free}.
\]
Its fully marginalized likelihood is
\[
    p(\data\mid\Theta)
    =
    \sum_{\bm{z}\in\{0,1\}^{N_{\rm cand}}}
    p(\data\mid\Theta,\bm{z})
    \prod_{i=1}^{N_{\rm cand}}
    q_i^{z_i}(1-q_i)^{1-z_i}.
\]
This is the direct analog of a two-component DFM-style mixture model, but here
each latent state changes a global lens model rather than an independent
single datum. The exact sum has $2^{N_{\rm cand}}$ global likelihood
evaluations and is not practical.

\subsection{Implemented one-candidate likelihood-ratio approximation}

The implementation expands around the all-scaling lens model. Let
\[
    \Theta_0
\]
denote the state in which every independent/free candidate is forced to its
scaling branch, i.e. $q_i$ is overridden to zero only for building this base
lens state. The base log likelihood is
\[
    \ell_0
    =
    \ell_{\rm core}(\Theta_0).
\]
For each candidate $i$, construct a one-free replacement state in which only
candidate $i$ is replaced by its free branch:
\[
    \bm{\alpha}^{(i)}
    =
    \bm{\alpha}^{(0)}
    -
    \bm{\alpha}_{S_i}
    +
    \bm{\alpha}_{F_i}.
\]
In local-Jacobian mode the same replacement is applied to the Jacobian field:
\[
    A^{(i)}
    =
    A^{(0)}
    -
    J_{S_i}
    +
    J_{F_i}.
\]
In stage 2 source-plane mode, the replacement is only the source-coordinate
replacement induced by the deflection-field change. There is no local-Jacobian
replacement term in the stage-2 likelihood.
The one-free log likelihood is
\[
    \ell_i^{F}
    =
    \ell_{\rm core}(\Theta^{(i)}).
\]
The per-candidate likelihood ratio is
\[
    \Delta_i
    =
    \ell_i^{F}-\ell_0.
\]
For numerical stability the implementation clips $\Delta_i$ to
$[-\Delta_{\max},\Delta_{\max}]$ before applying the mixture formula.

The population-mode approximation is
\[
    \boxed{
    \log p(\data\mid\Theta,{\rm population})
    =
    \ell_0
    +
    \sum_{i=1}^{N_{\rm cand}}
    \log\left[
        (1-q_i)
        +
        q_i \exp(\Delta_i)
    \right].
    }
\]
Equivalently, the code evaluates each summand stably as
\[
    \operatorname{logaddexp}
    \left[
        \log(1-q_i),
        \log(q_i)+\Delta_i
    \right].
\]

This is exact if the joint likelihood ratio factorizes as
\[
    \log p(\data\mid\Theta,\bm{z})-\ell_0
    =
    \sum_i z_i\Delta_i.
\]
In the real lens model this factorization is approximate, because two galaxies
changed at the same time can interact through the common source centroids and
the common lens field. The approximation keeps the closest tractable analog of
the DFM mixture: it compares each candidate's scaling branch to its free branch
using the real global likelihood, but only one candidate at a time.

\subsection{Posterior free-branch membership}

Population mode also has an explicit posterior free-branch membership
diagnostic:
\[
    r_i
    =
    p(z_i=1\mid\data,\Theta)
    \approx
    \frac{
        q_i\exp(\Delta_i)
    }{
        (1-q_i)+q_i\exp(\Delta_i)
    }.
\]
This $r_i$ is the useful classification quantity in population mode. The raw
sigmoid $q_i$ remains the prior branch weight implied by the population gate
parameters at that posterior sample.

\section{Priors and full sampled target}

The sampler target adds parameter priors to the likelihood above:
\[
    \log \pi(\Theta\mid\data)
    =
    \log p(\data\mid\Theta)
    +
    \sum_k \log \pi_k(\Theta_k)
    +
    {\rm const}.
\]
This includes the priors on the gate intercepts, magnitude slope, local logit
offsets, and the free-branch dPIE parameters. The distinction is that in blend
mode $q_i$ affects the likelihood only by changing the continuous lens field,
whereas in population mode $q_i$ also appears explicitly in the marginalized
mixture term.

\section{Why blend is fast}

Blend mode is fast because it performs one lens-state construction and one
core likelihood evaluation per parameter vector:
\[
    \Theta
    \longrightarrow
    \{\bm{\beta}_j,A_j\}_{j=1}^{N_{\rm img}}
    \longrightarrow
    \ell_{\rm core}.
\]
All independent/free candidates are folded into the same vectorized lens field
through their continuous branch weights. Computationally, the expensive part is
paid once.

With the refreshing surrogate path, cached inactive contributions and active
exact contributions are also combined once. There is no candidate loop over
alternative global likelihoods.

\section{Why population is slow}

Population mode is slower because it must estimate a likelihood ratio for each
candidate:
\[
    \Delta_i =
    \ell_{\rm core}(\text{candidate }i\text{ free, all others scaling})
    -
    \ell_{\rm core}(\text{all scaling}).
\]
Thus one parameter vector requires:
\begin{itemize}
    \item one all-scaling core likelihood evaluation;
    \item construction of one-free replacement source coordinates for all
    candidates;
    \item one additional core likelihood evaluation per candidate, vectorized
    with \texttt{jax.vmap} but still doing work proportional to
    $N_{\rm cand}$;
    \item in local-Jacobian mode, replacement and evaluation of the Jacobian
    covariance terms for each candidate as well.
\end{itemize}
The operation is vectorized rather than written as a Python loop, but the
amount of numerical work still scales roughly like
\[
    O(N_{\rm cand}N_{\rm img})
\]
for the replacement likelihoods, on top of the base likelihood. Local-Jacobian
mode has larger constants because each candidate replacement carries four
Jacobian arrays and full two-dimensional covariance calculations.

The stage-2 source-plane population mode is the cheaper version of the same
idea. It still evaluates one replacement likelihood per candidate, so it is
slower than blend mode, but each replacement only needs the source-plane
centroid and diagonal source covariance calculation. Stage-3 local-Jacobian
population mode does the same candidate-by-candidate comparison while also
updating $A_j$ and recomputing the local covariance matrices $C_j$, so it is
substantially more expensive.

\section{Summary}

\begin{center}
\begin{tabular}{lll}
\toprule
Stage/mode & Likelihood form & Interpretation of $q_i$\\
\midrule
stage 2 blend &
$\ell_{\rm src}$ of a continuously blended lens field &
continuous branch weight\\
stage 2 population &
$\ell_0^{\rm src}+\sum_i\log[(1-q_i)+q_i e^{\Delta_i^{\rm src}}]$ &
prior free-branch probability\\
stage 3 blend &
$\ell_{\rm core}$ of a continuously blended lens field &
continuous branch weight\\
stage 3 population &
$\ell_0+\sum_i\log[(1-q_i)+q_i e^{\Delta_i}]$ &
prior free-branch probability\\
\bottomrule
\end{tabular}
\end{center}

Blend is fast because it evaluates one blended global lens model. Population is
slower because it evaluates an approximate one-candidate free/scaling
likelihood ratio for every candidate. Stage-2 source-plane population is the
fastest population variant because it uses only $\ell_{\rm src}$. Stage-3
local-Jacobian population is slower because it also propagates each candidate
replacement through the local covariance. Population is closer to the desired
latent branch-membership mixture model; blend is a cheaper continuous
relaxation.

\end{document}
